% !TeX root = ../../Book4.tex
% 纲要标题：### 6.2 复形的运算
% 纲要标签：b4:sec:0502
% 纲要位置：计划纲要.md，第 3224 行。

\section{复形的运算}
\label{b4:sec:0502}

复形之间的态射必须与微分相容，Hom 和张量积的微分还必须消去两个方向的混合项。逐次的核与余核可以沿用交换范畴的构造，分次符号却需要另作核验。

% 纲要内容（仅作写作计划，尚非教材正文）：
%
% * 复形映射、核、余核与短正合列
% * Hom 复形及张量复形的初步定义
% * 分次符号规则
%

\subsection{复形映射、核、余核与短正合列}
\label{b4:subsec:0502:01}

\begin{definition}\label{b4:def:complex-map}
设 \(C,D\) 是交换范畴 \(\mathcal A\) 中的上链复形。态射族 \(f^n:C^n\to D^n\) 若满足
\[
\mathrm{d}_D^nf^n=f^{n+1}\mathrm{d}_C^n\qquad(n\in\mathbb Z),
\]
则称为\term[fuxing-yingshe]{复形映射}[map of complexes]，记为 \(f:C\to D\)。以逐次复合作为复合，得到复形范畴 \(\operatorname{Ch}(\mathcal A)\)。
\end{definition}
\symidx{\(\operatorname{Ch}(\mathcal A)\)：上链复形及复形映射组成的范畴}

记余圈嵌入为 \(z_X^n:Z^n(X)\to X^n\)（\(X=C,D\)）。交换等式给出 \(\mathrm{d}_D^nf^nz_C^n=f^{n+1}\mathrm{d}_C^nz_C^n=0\)，所以核的泛性质唯一产生
\[
Z^n(f):Z^n(C)\longrightarrow Z^n(D),
\qquad z_D^nZ^n(f)=f^nz_C^n.
\]
另一方面，\(f^n\mathrm{d}_C^{n-1}=\mathrm{d}_D^{n-1}f^{n-1}\) 保证 \(f^n\) 将前一微分的像送入目标的余边对象，因而 \(Z^n(f)\) 将 \(B^n(C)\) 的嵌入送入 \(B^n(D)\)。记 \(q_X^n:Z^n(X)\to H^n(X)\) 为余核态射，则 \(q_D^nZ^n(f)\) 杀死 \(B^n(C)\) 的嵌入，沿 \(q_C^n\) 唯一下降为
\[
H^n(f):H^n(C)\longrightarrow H^n(D),
\qquad H^n(f)q_C^n=q_D^nZ^n(f).
\]
在模范畴中，这两次因子化正是先将余圈 \(z\) 送到 \(f^n(z)\)，再取类 \([z]\mapsto[f^n(z)]\)。恒等映射和复合的保持也由这两种唯一性得到，因此 \(H^n\) 是函子。

\begin{proposition}\label{b4:prop:complex-category-abelian}
若 \(\mathcal A\) 是交换范畴，则 \(\operatorname{Ch}(\mathcal A)\) 也是交换范畴。核、余核、像与有限双积均逐次数计算；特别地，复形序列 \(0\to A\to B\to C\to0\) 短正合，当且仅当每个次数的对象序列短正合。
\end{proposition}
\begin{proof}
给定复形 \(C,D\)，有
\[
\operatorname{Hom}_{\operatorname{Ch}(\mathcal A)}(C,D)
\subseteq\prod_{n\in\mathbb Z}\operatorname{Hom}_{\mathcal A}(C^n,D^n).
\]
由于 \(\mathcal A\) 局部小，右端是集合指标的集合直积，故左端也是集合。这证明了 \(\operatorname{Ch}(\mathcal A)\) 的局部小性。态射逐次相加，零复形和逐次双积给出加性结构。

对 \(f:C\to D\)，令 \(K^n=\ker f^n\)，核嵌入为 \(k^n:K^n\to C^n\)。微分 \(\mathrm{d}_C^n\) 接在这个嵌入之后，得到从 \(K^n\) 到 \(C^{n+1}\) 的态射，且
\[
f^{n+1}\mathrm{d}_C^nk^n=\mathrm{d}_D^nf^nk^n=0.
\]
因此这条态射唯一地通过 \(f^{n+1}\) 的核分解，产生 \(\mathrm{d}_K^n:K^n\to K^{n+1}\)，满足 \(k^{n+1}\mathrm{d}_K^n=\mathrm{d}_C^nk^n\)。再计算
\[
k^{n+2}\mathrm{d}_K^{n+1}\mathrm{d}_K^n
=\mathrm{d}_C^{n+1}\mathrm{d}_C^nk^n=0,
\]
由 \(k^{n+2}\) 单，得到 \(\mathrm{d}_K^{n+1}\mathrm{d}_K^n=0\)。若复形映射 \(h:T\to C\) 满足 \(fh=0\)，其逐次核分解在复合 \(k^{n+1}\) 后由 \(h\) 的微分相容性互相匹配；消去这个单态射，便知分解本身是复形映射。故所得 \(K\to C\) 满足复形范畴中的核泛性质。

对余核 \(q^n:D^n\to Q^n\)，有 \(q^{n+1}\mathrm{d}_D^nf^n=q^{n+1}f^{n+1}\mathrm{d}_C^n=0\)。所以 \(q^{n+1}\mathrm{d}_D^n:D^n\to Q^{n+1}\) 杀死 \(f^n\)，沿其余核 \(q^n\) 唯一下降，产生 \(\mathrm{d}_Q^n:Q^n\to Q^{n+1}\)，满足
\[
\mathrm{d}_Q^nq^n=q^{n+1}\mathrm{d}_D^n.
\]
于是 \(\mathrm{d}_Q^{n+1}\mathrm{d}_Q^nq^n=q^{n+2}\mathrm{d}_D^{n+1}\mathrm{d}_D^n=0\)，消去满态射 \(q^n\) 得到复形条件。任何杀死 \(f\) 的复形映射都逐次通过 \(q^n\) 唯一分解；所得分解的微分相容性也在复合 \(q^n\) 后成立，再由其满性得到。因此 \(D\to Q\) 是复形范畴中的余核。

像是余核的核，余像是核的余核，所以二者也逐次计算。记 \(P=\operatorname{coim}f\)、\(J=\operatorname{im}f\)，结构复形映射为 \(p:C\to P\)、\(i:J\to D\)。每个次数的典范比较 \(a^n:P^n\to J^n\) 满足 \(i^na^np^n=f^n\)，且因 \(\mathcal A\) 交换而可逆。由 \(p,i,f\) 的微分相容性，得到
\[
\begin{aligned}
i^{n+1}\mathrm{d}_J^na^np^n
&=\mathrm{d}_D^nf^n=f^{n+1}\mathrm{d}_C^n\\
&=i^{n+1}a^{n+1}\mathrm{d}_P^np^n.
\end{aligned}
\]
消去单态射 \(i^{n+1}\) 与满态射 \(p^n\)，可知 \(\mathrm{d}_J^na^n=a^{n+1}\mathrm{d}_P^n\)，故这些比较组成复形映射。等式两侧再复合相应的逆，得到
\[
\mathrm{d}_P^n(a^n)^{-1}=(a^{n+1})^{-1}\mathrm{d}_J^n.
\]
所以逐次的逆也组成复形映射，比较确为复形同构。最后的正合性判据来自逐次核和像的识别。
\end{proof}

后面讨论长正合列与导出函子，还需要同调保持 Hom 群的加法。对复形映射 \(f,g:C\to D\)，记它们在余圈对象上诱导的态射为 \(Z^n(f),Z^n(g)\)。核的唯一分解给出 \(Z^n(f+g)=Z^n(f)+Z^n(g)\)。记余核态射为 \(q_X^n:Z^n(X)\to H^n(X)\)（\(X=C,D\)），则
\[
\begin{aligned}
H^n(f+g)q_C^n
&=q_D^n\bigl(Z^n(f)+Z^n(g)\bigr)\\
&=\bigl(H^n(f)+H^n(g)\bigr)q_C^n.
\end{aligned}
\]
由于 \(q_C^n\) 满，得到 \(H^n(f+g)=H^n(f)+H^n(g)\)。所以每个 \(H^n:\operatorname{Ch}(\mathcal A)\to\mathcal A\) 都是加性函子。

\ProseParagraphBreak
逐次分裂的短正合列不必有复形截面。微分可能把各次数所选补对象送出补对象；\secref{b4:sec:0503}中的非零连接同态能够检测这类不分裂现象。不过，连接同态全部为零一般也不足以推出存在复形截面。

\subsection{Hom 复形与张量复形}
\label{b4:subsec:0502:02}

给两个复形逐项取 Hom 或张量积，并不能只把所有微分原样相加：两个方向的混合项必须抵消。以下分别固定产生这种抵消的符号。

\ProseParagraphBreak
次数增加 \(m\) 的分次同态需要同时给出每个次数上的分量 \(f^p:C^p\to D^{p+m}\)，不要求只有有限多个分量非零，因此这些同态组成直积。例如取每个 \(p\in\mathbb Z\) 处都有 \(C^p=\mathbb Z\)，并令微分为零，则恒等映射 \(1_C=(1_{C^p})_{p\in\mathbb Z}\) 有无穷多个非零分量。它属于 \(\prod_p\operatorname{Hom}_{\mathbb Z}(C^p,C^p)\)，却不属于相应直和。

\ProseParagraphBreak
Hom 微分应把次数 \(m\) 的分次同态送到次数 \(m+1\)，而在次数零，微分为零应当恰好表示 \(\mathrm{d}_Df=f\mathrm{d}_C\)。因此取目标微分后的复合与源微分前的复合之差；为了让两次微分中的混合项抵消，第二项的系数随次数增加而变号，得到下面的带符号公式。

\begin{definition}\label{b4:def:hom-complex}
设 \(R\) 是含幺环，\(C,D\) 是左 \(R\)-模的上链复形。定义交换群复形
\[
\operatorname{Hom}_R^m(C,D)\coloneqq
\prod_{p\in\mathbb Z}\operatorname{Hom}_R(C^p,D^{p+m}),
\qquad
(\mathrm{d}f)^p=\mathrm{d}_D^{p+m}f^p-(-1)^m f^{p+1}\mathrm{d}_C^p.
\]
它称为\term[Hom-fuxing]{Hom 复形}[Hom complex]。这里 \(m\in\mathbb Z\)，\(f=(f^p)_{p\in\mathbb Z}\in\operatorname{Hom}_R^m(C,D)\) 是一整族使次数增加 \(m\) 的同态，不要求其与微分交换。
\end{definition}
\symidx{\(\operatorname{Hom}_R^\bullet(C,D),\ \operatorname{Hom}_R^m(C,D)\)：Hom 复形及其第 \(m\) 次项}

\(\mathrm{d}f\) 的次数已是 \(m+1\)，所以再次施行微分时使用的符号是 \((-1)^{m+1}\)。暂略各分量的指标，展开为
\[
\begin{aligned}
\mathrm{d}(\mathrm{d}f)
={}&\mathrm{d}_D^2f-(-1)^m\cdot\mathrm{d}_Df\mathrm{d}_C\\
&-(-1)^{m+1}\cdot\mathrm{d}_Df\mathrm{d}_C
+(-1)^{2m+1}f\mathrm{d}_C^2=0.
\end{aligned}
\]
次数零的余圈正是复形映射。若 \(\mathcal A\) 为一般局部小交换范畴，也可把上述 Hom 换为其态射交换群，得到同样的交换群复形。

\begin{definition}\label{b4:def:tensor-complex}
设 \(R\) 是含幺环，\(C\) 是右 \(R\)-模的上链复形，\(D\) 是左 \(R\)-模的上链复形。定义交换群复形
\[
\begin{gathered}
(C\otimes_R D)^n\coloneqq\bigoplus_{p+q=n}C^p\otimes_R D^q,\\
\mathrm{d}(x\otimes y)=\mathrm{d}_Cx\otimes y+(-1)^p x\otimes \mathrm{d}_Dy,\\
x\in C^p,\quad y\in D^q,\quad p,q\in\mathbb Z.
\end{gathered}
\]
它称为\term[zhangliang-fuxing]{张量复形}[tensor product of complexes]。若 \(R\) 交换且两者均为 \(R\)-模复形，结果也带自然的 \(R\)-模结构。
\end{definition}
\symidx{\(C\otimes_RD\)：复形的张量积}

每个张量项由整数对 \((p,q)\) 标记，取 \(p+q=n\) 的全部项便是总次数 \(n\)。\(\mathrm{d}_C\) 将 \(p\) 增一，\(\mathrm{d}_D\) 将 \(q\) 增一，所以二者都将总次数增一。虽然同一个总次数中可能有无限多个非零项，直和的定义只允许每个元素使用其中有限多个；微分后仍只有有限多个分量。公式与张量积的平衡关系相容，因为微分均为 \(R\)-线性。

\ProseParagraphBreak
混合项有两种产生次序：先作用 \(\mathrm{d}_C\)，则 \(x\) 的次数已经变为 \(p+1\)，接着作用 \(\mathrm{d}_D\) 时带系数 \((-1)^{p+1}\)；反过来先作用 \(\mathrm{d}_D\)，带系数 \((-1)^p\)，再作用 \(\mathrm{d}_C\)。因此它们相消，完整计算为
\[
\mathrm{d}^2(x\otimes y)=\mathrm{d}_C^2x\otimes y
+\bigl((-1)^{p+1}+(-1)^p\bigr)\cdot\mathrm{d}_Cx\otimes \mathrm{d}_Dy
+x\otimes \mathrm{d}_D^2y=0.
\]
这证明了复形条件。张量与同调能否交换是另一问题：后面的\cref{b4:thm:kunneth-pid}将表明，在主理想整环上即使复形各项自由，同调的张量积仍可能漏掉 Tor 项。

\subsection{分次符号规则}
\label{b4:subsec:0502:03}

在本书所用的分次张量结构中，一个次数为 \(r\) 的齐次因子跨过次数为 \(s\) 的齐次因子时，带系数 \((-1)^{rs}\)；相应齐次映射的张量运算也遵守这一约定。这称为\term[fenci-fuhao-guize]{分次符号规则}[graded sign rule]。它适用于具体的交换或张量操作，不能只因公式中出现两个次数就任意加号；是否与微分相容，仍须从所定义的映射核验。

\begin{proposition}\label{b4:prop:graded-leibniz}
设 \(R\) 为含幺环，\(C,D,E\) 是幺性左 \(R\) 模的上链复形，\(r,s\in\mathbb Z\)，\(f\in\operatorname{Hom}_R^r(C,D)\)、\(g\in\operatorname{Hom}_R^s(D,E)\)。以 \((gf)^p=g^{p+r}f^p\)（\(p\in\mathbb Z\)）定义复合，则
\[
\mathrm{d}(gf)=(\mathrm{d}g)f+(-1)^s g(\mathrm{d}f).
\]
这里 \(gf=g\circ f\)，即先作用 \(f\)，再作用 \(g\)；因子在式中仍依次写作 \(g,f\)。在通常的分次乘法公式 \(\mathrm{d}(ab)=(\mathrm{d}a)b+(-1)^{|a|}a\cdot\mathrm{d}b\) 中取 \(a=g\)、\(b=f\)，左因子的次数便是 \(s\)，所以这里出现 \((-1)^s\)。
若 \(R\) 交换且 \(C,D\) 为 \(R\)-模复形，则
\[
\tau:C\otimes_R D\longrightarrow D\otimes_R C,
\qquad \tau(x\otimes y)=(-1)^{pq}y\otimes x
\quad(x\in C^p,\ y\in D^q)
\]
是复形同构，且 \(\tau_{D,C}\tau_{C,D}=1\)。
\end{proposition}
\begin{proof}
第一式右端展开为
\[
\begin{aligned}
(\mathrm{d}g)f+(-1)^s g(\mathrm{d}f)
={}&(\mathrm{d}_Eg-(-1)^sg\mathrm{d}_D)f\\
&+(-1)^sg(\mathrm{d}_Df-(-1)^rf\mathrm{d}_C)\\
={}&\mathrm{d}_Egf-(-1)^{r+s}gf\mathrm{d}_C.
\end{aligned}
\]
这正是左端。第二式的两种复合在纯张量上分别为
\[
\begin{aligned}
\tau\mathrm{d}(x\otimes y)
={}&(-1)^{(p+1)q}y\otimes\mathrm{d}_Cx\\
&+(-1)^{p+p(q+1)}\cdot\mathrm{d}_Dy\otimes x,\\
\mathrm{d}\tau(x\otimes y)
={}&(-1)^{pq}\cdot\mathrm{d}_Dy\otimes x\\
&+(-1)^{pq+q}y\otimes\mathrm{d}_Cx.
\end{aligned}
\]
\(y\otimes\mathrm{d}_Cx\) 的两份系数相等，因为 \((p+1)q=pq+q\)；\(\mathrm{d}_Dy\otimes x\) 的两份系数相等，因为 \(p+p(q+1)=pq+2p\)。因此 \(\tau\) 与微分相容。再交换一次乘上 \((-1)^{2pq}=1\)，所以映射可逆。
\end{proof}

例如两个只在次数一非零的 \(k\)-向量空间复形交换时，交换映射带负号。若删掉该号，在存在非零微分的复形上就可能不再与微分交换。特征二时正负号数值相同，但定义仍沿用同一公式，以便不同特征的论证可以比较。
